\section{Trabalhos Relacionados}
\label{sec:relacionados}

Os trabalhos relacionados se dividem em três grupos: jogos para o ensino de teste de software, jogos
educacionais desenvolvidos e avaliados em trabalhos de conclusão de curso da UFPB, e o uso de agentes
artificiais para avaliar e balancear jogos.

\subsection{Jogos para o ensino de teste de software}

O \emph{Debug} é um jogo de cartas cooperativo, criado numa disciplina de graduação, para explorar níveis e
tipos de teste de software \cite{costa2024debug}. O \emph{GAMUT} usa uma narrativa de fantasia medieval para
ensinar teste de unidade \cite{gomes2021gamut}. No \emph{GreaTest}, os jogadores assumem o papel de analistas de
teste e escolhem o tipo de teste adequado a cada cenário \cite{beppe2018greatest}. O \emph{Sojourner under
Sabotage} é um jogo sério no navegador em que o jogador corrige componentes sabotados de uma nave usando teste
de unidade e depuração \cite{straubinger2025sojourner}. O \emph{JADE} é um jogo de tabuleiro para ensinar
ergonomia de software \cite{jeandaubias2023jade}. Experiências com gamificação contínua em disciplinas de teste
também mostram ganho de engajamento \cite{straubinger2024ci,jesus2020,blanco2023}.

A primeira versão do Critical Failure \cite{oliveira2025critical} ocupou uma lacuna desses trabalhos, voltados
principalmente ao teste funcional: é um jogo de cartas competitivo sobre teste não funcional e atributos de
qualidade. Os jogos citados foram avaliados pela percepção dos estudantes. Nenhum deles relata uma verificação
sistemática do equilíbrio entre os jogadores antes da aplicação em sala.

\subsection{Jogos educacionais em trabalhos de conclusão de curso da UFPB}

No repositório institucional da UFPB há trabalhos de conclusão de curso que desenvolvem jogos de cartas para o
ensino de computação: o \emph{Computer Cards}, sobre arquitetura de computadores no ensino fundamental
\cite{oliveira2022computercards}; o \emph{BI-NA-RI-OH}, sobre aritmética binária \cite{soares2023binarioh}; e o
\emph{Logic Challenge}, sobre lógica \cite{abreu2021logic}. Outro trabalho avalia um jogo de tabuleiro para o
ensino de redes de computadores com o modelo MEEGA+ \cite{martins2020redes,petri2017meega}. No Centro de
Informática, um trabalho aplicou testes automatizados a um jogo sério, o GiroJampa, com testes de unidade e
medição de cobertura \cite{ramos2023girojampa}. Nenhum desses trabalhos trata de jogo para o ensino de teste de
software nem de balanceamento verificado por simulação.

\subsection{Balanceamento e playtesting com agentes}

O uso de agentes artificiais para avaliar o projeto de jogos tem uma linha de pesquisa própria.
Jaffe et al., por exemplo, mediram o balanceamento de jogos competitivos comparando a força de
jogadores com restrições impostas, numa abordagem chamada \emph{restricted play} \cite{jaffe2012}.
Silva et al. usaram agentes com heurísticas diferentes para jogar grandes quantidades de partidas
de um jogo de tabuleiro moderno e encontrar problemas de projeto, como estratégias que dominam e situações não
previstas pelas regras \cite{silva2017ttr}. Volz et al. mostraram que é possível balancear
automaticamente o baralho de um jogo de cartas com algoritmos evolutivos e simulação \cite{volz2016}.
Holmgård et al. propuseram \emph{personas} procedurais, agentes com estilos de jogo distintos, para
testar automaticamente os níveis de um jogo \cite{holmgard2019}.

Este trabalho se apoia nessa linha, com três diferenças. Primeiro, o jogo é educacional, e parte das decisões do
jogador depende de conhecimento sobre o conteúdo, como explicar um teste ou diagnosticar um incidente. Por isso
o simulador modela o conhecimento do estudante como probabilidade de acerto. Segundo, o balanceamento é
escrito como requisitos com critério de aceite fixado antes da medição, como em um teste de aceitação. Terceiro,
o próprio simulador é verificado com técnicas de teste de software, como testes metamórficos, testes baseados
em propriedades e teste de mutação. A Tabela~\ref{tab:relacionados} resume a comparação.

\begin{table}[ht]
  \centering
  \caption{Comparação entre os trabalhos relacionados e este trabalho}
  \label{tab:relacionados}
  \small
  \begin{tabularx}{\linewidth}{>{\raggedright\arraybackslash}p{3.9cm} >{\raggedright\arraybackslash}p{2.2cm}
      >{\raggedright\arraybackslash}p{1.9cm} >{\raggedright\arraybackslash}X >{\raggedright\arraybackslash}p{2.5cm}}
    \toprule
    Trabalho & Tipo & Modo & Escopo & Balanceamento verificado \\
    \midrule
    Debug \cite{costa2024debug} & Analógico (cartas) & Cooperativo & Níveis e tipos de teste & Não \\
    GAMUT \cite{gomes2021gamut} & Digital (narrativa web) & Individual & Teste de unidade & Não \\
    GreaTest \cite{beppe2018greatest} & Analógico (cartas) & Competitivo & Teste funcional em geral & Não \\
    Sojourner under Sabotage \cite{straubinger2025sojourner} & Digital (jogo sério) & Individual & Teste de unidade e depuração & Não \\
    JADE \cite{jeandaubias2023jade} & Analógico (tabuleiro) & Cooperativo & Ergonomia de software & Não \\
    Critical Failure v1 \cite{oliveira2025critical} & Analógico (cartas) & Competitivo & Teste não funcional e atributos de qualidade & Por observação (6 partidas) \\
    Agentes para balanceamento \cite{jaffe2012} \cite{silva2017ttr} \cite{volz2016} & Jogos de entretenimento & Competitivo & Projeto de jogos & Por simulação \\
    \textbf{Este trabalho} & Analógico (cartas) e simulador & Competitivo & Teste não funcional (ISO/IEC 25010:2023) & Por simulação, com requisitos e critério de aceite \\
    \bottomrule
  \end{tabularx}
\end{table}
